\documentclass[10pt,a4paper]{article}
\usepackage[T1]{fontenc}
\usepackage[utf8]{inputenc}
\usepackage[brazilian]{babel}
\usepackage[a4paper,margin=1.5cm]{geometry}
\usepackage{booktabs}
\usepackage{xcolor}
\usepackage{microtype}
\usepackage[hidelinks]{hyperref}
\pagestyle{empty}
\setlength{\parindent}{0pt}
\setlength{\parskip}{2pt}
% Paleta monocromática: só a escala gray do xcolor, para impressão em toner.
\definecolor{cinza}{gray}{0.35}
\definecolor{cinzaclaro}{gray}{0.94}
\definecolor{borda}{gray}{0.75}
\definecolor{fundocard}{gray}{0.98}
\definecolor{barras}{gray}{0.15}
\newcommand{\barra}[1]{\textcolor{barras}{\rule{#1mm}{2.4pt}}}
% Cartão de indicador: rótulo, número em destaque e rodapé de contexto.
\newcommand{\kpicard}[3]{%
  \fcolorbox{borda}{fundocard}{%
    \begin{minipage}[c][1.4cm][c]{\dimexpr0.32\textwidth-2\fboxsep-2\fboxrule\relax}
      \centering
      {\scriptsize\bfseries\color{cinza} #1}\\[2pt]
      {\Large\bfseries #2}\\[1.5pt]
      {\tiny\color{cinza} #3}%
    \end{minipage}%
  }%
}
\begin{document}

{\Large\bfseries Sete Barras --- climatologia das médias}\\[2pt]
{\footnotesize\color{cinza} Estação CIIAGRO/IAC ciiagro:80 \,\textbullet\, 321 meses (2000-03 a 2026-11)}\\[1pt]
{\footnotesize \textbf{25 anos com os 12 meses fechados} --- médias do período medido, não previsão do tempo.}

\vspace{4pt}
\colorbox{cinzaclaro}{%
  \parbox{\dimexpr\textwidth-2\fboxsep}{%
    \footnotesize
    \textbf{Leia antes de usar:} os números abaixo são \textbf{médias de vários anos}, obtidas das medições mensais do CIIAGRO/IAC. \textbf{Não é previsão meteorológica} para os próximos dias nem para esta safra.\par\vspace{2pt}
    \textbf{Aplicação prática:} este é o comportamento médio esperado da estação ---
    base técnica para o calendário agrícola, janelas de calagem, o preparo do solo e
    o manejo conservacionista.%
  }%
}
\vspace{4pt}
\noindent
\kpicard{Chuva anual média}{1.511 mm/ano}{em todo o ano}\hfill
\kpicard{Mês mais chuvoso}{266 mm}{em janeiro}\hfill
\kpicard{Mês mais seco}{56 mm}{em agosto}

\vspace{2pt}
\noindent
\kpicard{Temperatura média}{23,2 \textdegree{}C}{média dos 12 meses}\hfill
\kpicard{Mês mais quente}{27,2 \textdegree{}C}{em fevereiro}\hfill
\kpicard{Mês mais frio}{19,0 \textdegree{}C}{em julho}

\vspace{4pt}
\begin{center}
{\bfseries\small Tabela 1: Chuva média por mês do ano}\\[2pt]
{\small
\begin{tabular}{@{}l r l r r r@{}}
\toprule
mês & chuva (mm) & perfil & t mín & t méd & t máx \\
\midrule
janeiro    &   266 & \barra{26.0}   & 21,6 & 26,7 & 31,9 \\
fevereiro  &   168 & \barra{16.4}   & 21,8 & 27,2 & 32,6 \\
março      &   168 & \barra{16.4}   & 21,0 & 26,1 & 31,2 \\
abril      &    95 & \barra{9.3}    & 19,4 & 24,4 & 29,3 \\
maio       &    86 & \barra{8.4}    & 16,3 & 21,2 & 26,1 \\
junho      &    75 & \barra{7.3}    & 15,1 & 19,8 & 24,6 \\
julho      &    84 & \barra{8.2}    & 13,9 & 19,0 & 24,1 \\
agosto     &    56 & \barra{5.5}    & 14,8 & 19,9 & 25,0 \\
setembro   &   104 & \barra{10.2}   & 16,6 & 21,2 & 25,8 \\
outubro    &   138 & \barra{13.5}   & 18,4 & 22,8 & 27,3 \\
novembro   &   125 & \barra{12.2}   & 19,3 & 24,3 & 29,4 \\
dezembro   &   161 & \barra{15.7}   & 20,8 & 26,1 & 31,4 \\
\midrule
ano & 1.511 & & & 23,2 & \\
\bottomrule
\end{tabular}}
\end{center}

\begin{center}
{\bfseries\small Tabela 2: Chuva por estação do ano}\\[2pt]
{\small
\begin{tabular}{@{}l l r r r l l@{}}
\toprule
estação & meses & total (mm) & \% do ano & média/mês & maior mês & menor mês \\
\midrule
Verão      & dez--fev &   595 &  39\% &  198 & 266 (jan) & 161 (dez) \\
Outono     & mar--mai &   350 &  23\% &  117 & 168 (mar) & 86 (mai) \\
Inverno    & jun--ago &   215 &  14\% &   72 & 84 (jul) & 56 (ago) \\
Primavera  & set--nov &   367 &  24\% &  122 & 138 (out) & 104 (set) \\
\bottomrule
\end{tabular}}
\end{center}

\vspace{1pt}
{\footnotesize\textbf{Antes dos solstícios} (acumulado desde janeiro):
inverno: janeiro a maio somam 784 mm (52\% do ano) \,\textbullet\, verão: junho a novembro somam 582 mm (39\% do ano). Dezembro fica depois do solstício de verão e entra só no total
anual.}

\vspace{2pt}
{\bfseries\small Como usar no planejamento nesta estação}

\footnotesize\begin{itemize}\setlength{\itemsep}{0.5pt}
  \item \textbf{Amostragem de solo e preparo mecânico}: Em abril, maio, junho, julho e agosto a chuva média fica entre 56 e 95 mm/mês \textemdash{} é a janela de menor encharcamento para retirar terra para laudo, arar e incorporar calcário com antecedência.
  \item \textbf{Calagem em cobertura (pasto e perene já formado)}: A partir de outubro, quando a média volta a passar de 120 mm/mês: aí há água para o corretivo descer no perfil sem enxurrada.
  \item \textbf{Plantio e transplante}: O verão (dezembro a fevereiro) concentra 595 mm, 39\% da chuva do ano \textemdash{} favorece o pegamento de mudas sem irrigação pesada; em áreas declivosas, reforce a conservação de solo nesse período.
  \item \textbf{Excesso de chuva e conservação}: janeiro passa de 200 mm em média (266 mm): canteiros altos, drenagem limpa e cobertura morta antes da temporada, para não perder solo e adubo no escoamento.
  \item \textbf{Calor, frio e fitossanidade}: fevereiro tem a maior média (27,2 \textdegree{}C) e julho a menor (19,0 \textdegree{}C, com mínima média de 13,9 \textdegree{}C). Antecipe o monitoramento de doenças fúngicas nos meses quentes e úmidos e proteja culturas sensíveis ao frio nos mais frios.
\end{itemize}

\vspace{1pt}
{\bfseries\small Vale para qualquer município da região}

\scriptsize\begin{itemize}\setlength{\itemsep}{0.5pt}
  \item \textbf{Média não é previsão}: Descrevem o padrão de vários anos, não o tempo deste ano: use para escolher a janela de cada operação, e decida no campo.
  \item \textbf{Corrija antes de precisar}: Calcário reage em meses; aplicado na janela seca, chega ao plantio já reagido.
  \item \textbf{Chuva em mm não é água no solo}: Textura, declive e cobertura morta decidem quanto do que cai chega à raiz.
  \item \textbf{Série curta é pista, não padrão}: Poucos anos fechados oscilam muito: compare com a estação vizinha antes de fixar calendário.
\end{itemize}

\vspace{1pt}
{\tiny\color{cinza}
Fonte: CIIAGRO/IAC (\href{http://ciiagro.org.br}{ciiagro.org.br}), série mensal;
chuva é o total do mês. Painel:
\href{https://climalitoralsul.netlify.app/\#sete-barras}%
{\texttt{climalitoralsul.netlify.app/\#sete-barras}}.
Umidade relativa omitida: a fonte usa \texttt{0} para ``sem sensor''. 
médias do período 2000\textendash{}2026 (26.8 anos de dado; parte da série não validada).
Gerado a partir de \texttt{dados.json} de 2026-11-06.}

\end{document}
